\chapter{Results and Discussion}
% TODO: no data gathered yet -- placeholder chapter.

``We know about `cults' largely by what the media tells us, and their views have been overwhelmingly negative\autocite{vandrielPrintMediaCoverage1988, pfeiferPsychologicalFramingCults1992}'' \autocite[pp.1]{dawsonCultsNewReligious2009}


\section{A priori observations}
Before any use of geometrical tools, simple statistic as well as qualitative observations of the study can be made.

\subsection{Survey}

\subsubsection{Sample}
20/26 (77\%) interviews were in person, 6/26 (23\%) via Instagram DM. 21/26 (81\%) are in English, 5/26 (19\%) in French — all French interviews fall in Batch 3 and have English translations in the database.

\subsubsection{Recurring themes}
\begin{itemize}
  \item \textbf{Charismatic/exploitative leader.} The single most consistent structural element across the corpus, in both languages.
  \item \textbf{Isolation and restricted freedom} as a recurring structural marker.
  \item \textbf{Cult vs.\ religion is unresolved.} Not a settled theme but a live argument interviewees have with themselves, often surfacing through comparisons to religion.
  \item \textbf{"Cult" as adjective.} The term has largely shifted to adjective use for devoted-but-harmless fandom. Non-native speakers map this onto \textit{Kult} (German) and \textit{culte} (French), applying it loosely to anything: bike parts, bands,that draws a passionate crowd.
  \item \textbf{Media as the main source of "cult" knowledge.} Documentaries and fiction (e.g.\ referenced Aug 6, 19:55) visibly shape how interviewees frame the concept.
\end{itemize}

\subsubsection{Divergences}
\begin{itemize}
  \item Positive vs.\ harm-centered valence of "cult" itself.
  \item Whether outward recruitment is necessary to qualify as a cult.
\end{itemize}

\subsubsection{Other notes}
Interviewees were often unsure what the term means precisely, and sarcastic use was rarer than expected. One (Aug 22, 21:01) used "dérive dangereuse," closely echoing the sectarian-drift concept. Unexpected recurring angles included privilege, Western perspective, and masculinity. Emergent topics align with the sectarian-drift literature: distortion of consent, a central leader figure, negative connotation, and use as an outside label rather than self-description.


\section{Global structure and cluster structure}
Before diving into the details of the results, it is useful to provide a high-level overview of the global structure of the data and the cluster structure that emerges from the analysis.

% insert pairwise centroid distance short table here, with a note that the full table is in the appendix.





We need to add the definitions to the glossary or in the methodology section. Here are some key terms used in the analysis:

Centroid: the mean vector of a set of points in the embedding space; the geometric centre of that set.

Pairwise (centroid-distance matrix): for every pair of centroids, the Euclidean distance between them; a symmetric table showing how far apart each pair of group centres is.

Euclidean distance: ordinary straight-line distance, the same idea as distance on a map,
just in many more than two dimensions. Smaller = more similar.

Cosine similarity: instead of measuring the gap between two points, this measures whether they point in the same direction from the origin, ignoring how long each vector is. It ranges from $-1$ (opposite directions) to $1$ (exactly the same direction), with 0 meaning unrelated/perpendicular directions.

k-nearest-neighbour composition
For a given point, its k nearest neighbours are simply the k closest other points in the
space (this document uses k = 10 and k = 20).

Silhouette score: A single number, in principle between $-1$ and $1$, that scores how well a proposed grouping
(here: the three expression sources) matches the actual geometric clustering of the data – roughly, "on average, is each point closer to others in its own group than to points in other groups?" A score near 0 means the proposed groups do not correspond to distinct geometric clusters; a score near 1 would mean they line up almost perfectly with the actual shape of the data






===

\section{Feedback on the results}
